In this section, we describe our evaluation framework\footnote{{https://github.com/haoyuhan1/RAGvsGraphRAG}}
 and experimental protocol. To ensure fair comparison, we evaluate RAG and GraphRAG under identical settings whenever applicable, and otherwise follow the default configurations of each method while matching key budgets. We decouple retrieval from generation by first saving retrieved evidence for each method and then using a unified generation script to produce outputs conditioned on the saved retrieval results.

\subsection{RAG Pipeline}
We adopt a standard dense-retrieval RAG pipeline~\cite{karpukhin2020dense}. 
Given a corpus, we segment documents into textual chunks and build an index by embedding each chunk into a shared vector space. 
At inference time, we embed the query, retrieve top-ranked chunks based on similarity.

\subsection{GraphRAG Implementations}
GraphRAG designs differ in how structures are constructed and how structural information is used during retrieval. 
In this work, we group GraphRAG approaches into four representative classes and select representative implementations for each class: 

\noindent \textbf{KG-based GraphRAG.}
In KG-based GraphRAG~\cite{Liu_LlamaIndex_2022}, a knowledge graph is constructed from text. 
Given a query, relevant entities are identified and aligned to nodes in the KG, and retrieval traverses multi-hop neighborhoods to collect relational triplets \textit{(head, relation, tail)} as evidence.
We consider two variants: 
{\bf (1)} {\it KG-GraphRAG (Triplets)}, which retrieves only triplets, and 
{\bf (2)} {\it KG-GraphRAG (Triplets+Text)}, which retrieves both triplets and their associated source text. 
We implement KG-GraphRAG using LlamaIndex~\cite{Liu_LlamaIndex_2022}.\footnote{{https://www.llamaindex.ai/}}

\noindent \textbf{Community-based GraphRAG.}
Community-based GraphRAG~\cite{edge2024local} further organizes the constructed KG into hierarchical communities using graph clustering algorithms. 
Each community is associated with a textual summary/report, where lower-level communities capture fine-grained information and higher-level communities provide increasingly abstract representations.
We evaluate two retrieval modes: 
{\bf Local Search} retrieves entity neighborhoods and lower-level community reports via entity matching, denoted as {\it Community-GraphRAG (Local)}; 
{\bf Global Search} retrieves high-level community summaries by semantic similarity, denoted as {\it Community-GraphRAG (Global)}.
We adopt the implementation of \citet{edge2024local}.\footnote{\url{https://microsoft.github.io/graphrag}}

\noindent \textbf{Text-centric graph-guided RAG.}
Text-centric graph-guided methods retain original text chunks as the primary retrieval units while leveraging graph structures to guide scoring or traversal. 
We select HippoRAG2~\cite{gutierrez2025rag} as a representative method. 
HippoRAG2 builds an entity-linked graph over chunks and retrieves query-relevant entities first, followed by the connected text chunks. 
Here, the graph acts as an auxiliary structure guiding chunk retrieval rather than the primary retrieval target.

\noindent \textbf{Hierarchical summary-based GraphRAG.}
Hierarchical summary-based methods construct multi-level hierarchical structures over text, where higher-level nodes represent progressively more abstract summaries of lower-level content. 
We adopt RAPTOR~\cite{sarthi2024raptor} as a representative method. 
RAPTOR recursively clusters text chunks and generates summaries at each level, enabling coarse-to-fine, multi-granular retrieval without relying on explicit KGs.

\subsection{Tasks}
We evaluate all methods on two representative text-based tasks: Question Answering and Query-based Summarization, covering both single-hop and multi-hop QA, as well as single-document and multi-document summarization. For single-document tasks, retrieval is restricted to the corresponding document, while for multi-document tasks, retrieval is performed over an index constructed from all documents.

\subsection{Unified Experimental Settings}
To ensure fair comparison across RAG and GraphRAG methods, we standardize core settings whenever applicable.

\noindent \textbf{Graph construction.}
For GraphRAG methods that require graph construction from text (e.g., KG-GraphRAG, Community-GraphRAG, and HippoRAG2), graphs are constructed using GPT-4o-mini; results with GPT-4o are reported in Appendix~\ref{app:graphconstruction}.

\noindent \textbf{Chunking.}
We segment documents into chunks of approximately \textit{256} tokens for all methods.

\noindent \textbf{Embedding model and retrieval budget.}
We embed queries, chunks, and graph information into a shared vector space using OpenAI's \texttt{text-embedding-ada-002} model~\cite{nussbaum2024nomic}, and retrieve top-$k$ candidates by semantic similarity, where $k{=}10$ by default. 

\noindent \textbf{Reranking.}
When reranking is enabled, we apply a cross-encoder reranker to reorder retrieved candidates and select the final top-$k$ evidence units. 
We use \texttt{BAAI/bge-reranker-large}~\cite{bge_embedding} as the reranker for all methods that support reranking, ensuring consistent reranking behavior across systems.

\noindent \textbf{Iterative retrieval.}
When iterative retrieval is enabled, we adopt IRCoT~\cite{trivedi2022interleaving}, which interleaves retrieval with intermediate reasoning steps. 

\noindent \textbf{Generation backbones.}
To control for generation capacity, we use two open-source instruction-tuned LLMs of different sizes as generators: Llama-3.1-8B-Instruct and Llama-3.1-70B-Instruct~\cite{dubey2024llama}. 
